% Part: first-order-logic
% Chapter: axiomatic-deduction
% Section: proof-theoretic-notions

\documentclass[../../../include/open-logic-section]{subfiles}

\begin{document}

\iftag{FOL}
      {\olfileid{fol}{axd}{ptn}}
      {\olfileid{pl}{axd}{ptn}}

\olsection{证明论概念}

\begin{explain}
正如我们已经定义了若干重要的语义概念
（\iftag{FOL}{有效性}{重言式}、语义后承、可满足性），我们现在定义
相对应的\emph{证明论概念}。这些概念不是诉诸语句在结构中的满足关系来定义的，
而是依据某些公式的\emph{可推导性}或\emph{不可推导性}来定义的。
一个重要的发现是，这些概念是重合的。这即是\emph{可靠性}与\emph{完备性}定理的内容。
\end{explain}

\begin{defn}[可推导性]
公式~$!A$ \emph{从 $\Gamma$ 可推导}，记作 $\Gamma \Proves !A$，
如果存在一个从~$\Gamma$ 出发并以~$!A$ 结尾的推导。
\end{defn}

\begin{defn}[定理]
公式~$!A$ 是一个\emph{定理}，如果存在一个从空集推导~$!A$ 的推导。
若~$!A$ 是定理，则记作 $\Proves !A$；否则记作 $\Proves/ !A$。
\end{defn}

\begin{defn}[一致性]
公式的集合~$\Gamma$ 是\emph{一致的}，当且仅当 $\Gamma\Proves/ \lfalse$，否则它是\emph{不一致的}。
\end{defn}

\begin{prop}[自反性]
\ollabel{prop:reflexivity}
若 $!A \in \Gamma$，则 $\Gamma \Proves !A$。
\end{prop}

\begin{proof}
  公式~$!A$ 本身即构成一个从~$\Gamma$ 推导~$!A$ 的推导。
\end{proof}

\begin{prop}[单调性]
\ollabel{prop:monotonicity}
若 $\Gamma \subseteq \Delta$ 且 $\Gamma \Proves !A$，则 $\Delta \Proves !A$。
\end{prop}

\begin{proof}
  任何从~$\Gamma$ 推导~$!A$ 的推导，同时也是从~$\Delta$ 推导~$!A$ 的推导。
\end{proof}

\begin{prop}[传递性]
\ollabel{prop:transitivity}
若 $\Gamma \Proves !A$ 且 $\{!A\} \cup \Delta \Proves
!B$，则 $\Gamma \cup \Delta \Proves !B$。
\end{prop}

\begin{proof}
  假设 $\{!A\} \cup \Delta \Proves !B$。那么存在一个从集合~$\{!A\} \cup \Delta$ 推导~$!B$ 的推导 $!B_1$, \dots, $!B_l = !B$。
  在该推导中，某些步骤的正确性来自于某条引用了前述某行~$!B_i = !A$ 的规则。
  根据假设，存在一个从~$\Gamma$ 推导~$!A$ 的推导，即一个推导~$!A_1$, \dots, $!A_k = !A$，其中每个 $!A_i$ 或是公理，或是~$\Gamma$ 中的元素，或是根据推理规则而正确的。
  现在考虑序列
  \[
  !A_1, \dots, !A_k = !A, !B_1, \dots, !B_l = !B.
  \]
  这是一个从 $\Gamma \cup \Delta$ 推导~$!B$ 的正确推导，因为现在每一个等于 $!A$ 的 $B_i$ 都由证成 $!A_k = !A$ 的同一条规则所证成。
\end{proof}

注意，这特别意味着，如果 $\Gamma \Proves !A$ 且 $!A
\Proves !B$，那么 $\Gamma \Proves !B$。由此也可得，如果 $!A_1,
\dots, !A_n \Proves !B$ 且对每个~$i$ 都有 $\Gamma \Proves !A_i$，那么
$\Gamma \Proves !B$。

\begin{prop}
\ollabel{prop:incons}
$\Gamma$ 是不一致的，当且仅当对每个~$!A$ 都有 $\Gamma \Proves !A$。
\end{prop}

\begin{proof}
习题。
\end{proof}

\tagprob{FOL}


\begin{prob}
证明 \olref[fol][axd][ptn]{prop:incons}。
\end{prob}


\tagendprob

\tagprob{notFOL}


\begin{prob}
证明 \olref[pl][axd][ptn]{prop:incons}。
\end{prob}


\tagendprob

\begin{prop}[紧致性]
\ollabel{prop:proves-compact}
  \begin{enumerate}
  \item 若 $\Gamma \Proves !A$，则存在有限子集 $\Gamma_0
    \subseteq \Gamma$ 使得 $\Gamma_0 \Proves !A$。
  \item 若 $\Gamma$ 的每个有限子集都是一致的，则 $\Gamma$ 是一致的。
  \end{enumerate}
\end{prop}

\begin{proof}
  \begin{enumerate}
    \item 若 $\Gamma \Proves !A$，则存在一个有限的公式序列 $!A_1$, \dots,~$!A_n$，使得 $!A \ident !A_n$，且每个 $!A_i$ 或是逻辑公理，或是~$\Gamma$ 中的元素，或是由先前的公式经肯定前件律得出。取 $\Gamma_0$ 为属于~$\Gamma$ 的那些 $!A_i$。那么该推导同样是从~$\Gamma_0$ 出发的推导，因此 $\Gamma_0 \Proves !A$。
    \item 这是 (1) 在 $!A \ident \lfalse$ 这一特殊情况下的逆否命题。
\end{enumerate}
\end{proof}

\end{document}